\begin{figure}[t]
\vspace{-3ex}
\begin{lstlisting}[
multicols=2,
xleftmargin=0em, 
numbers=none,
language=Python, 
]
class TTT_Layer(nn.Module):
  def __init__(self):
    self.task = Task()

  def forward(self, in_seq):
    state = Learner(self.task)
    out_seq = []
    for tok in in_seq:
      state.train(tok)
      out_seq.append(state.predict(tok))
    return out_seq

class Task(nn.Module):
  def __init__(self):
    self.theta_K = nn.Param((d1, d2))
    self.theta_V = nn.Param((d1, d2))
    self.theta_Q = nn.Param((d1, d2))

  def loss(self, f, x):
    train_view = self.theta_K @ x
    label_view = self.theta_V @ x
    return MSE(f(train_view), label_view)
class Learner():
  def __init__(self, task):
    self.task = task
    # Linear here, but can be any model
    self.model = Linear()
    # online GD here for simplicity
    self.optim = OGD()

  def train(self, x):
    # grad function wrt first arg
    # of loss, which is self.model
    grad_fn = grad(self.task.loss)    
    # calculate inner-loop grad
    grad_in = grad_fn(self.model, x)

    # starting from current params,
    # step in direction of grad_in, 
    self.optim.step(self.model, grad_in)
    
  def predict(self, x):
    test_view = self.task.theta_Q @ x
    return self.model(test_view)
\end{lstlisting}
\vspace{1.5ex}
\caption{
Naive implementation of a TTT layer with a linear model and online GD in the style of PyTorch.
\texttt{TTT\_Layer} can be dropped into a larger network like other sequence modeling layers. 
Training the network will optimize the parameters of \texttt{Task} in \texttt{TTT\_Layer}, because both are subclasses of \texttt{nn.Module}.
Since \texttt{Learner} is not a subclass of \texttt{nn.Module}, \texttt{state.model} is updated manually in the inner loop for each call of \texttt{state.train}.
For simplicity, we sometimes overload \texttt{model} as \texttt{model.parameters}.
}
\label{fig:code}
\end{figure}
